% ============================================================
% Capítulo 1 — Introducción
% ============================================================

\chapter{Introducción}
\label{cap:introduccion}

% ------------------------------------------------------------
\section{Motivación}
% ------------------------------------------------------------

Las redes neuronales de segmentación se entrenan casi siempre con
pérdidas de solapamiento, como la entropía cruzada o el coeficiente de
Dice, que comparan la predicción con la referencia píxel a píxel. Para
muchos objetos esto basta. Pero cuando lo que se segmenta son
estructuras finas y alargadas (vasos sanguíneos, carreteras, fibras),
los errores que más importan no son los que más píxeles cuestan. Un
vaso partido en dos por un hueco de tres píxeles apenas cambia el Dice,
y sin embargo ha dejado de ser un vaso para convertirse en dos. Son
errores de \emph{forma}: cambian cuántas piezas hay y cómo se conectan.

La topología es la parte de las matemáticas que estudia precisamente
esas propiedades, las que no cambian al deformar un objeto sin
romperlo ni pegarlo. En los últimos años, la topología computacional, y
en particular la homología persistente, ha dado lugar a métricas y
pérdidas que penalizan directamente los errores topológicos de una
segmentación~\cite{hu2019,shit2021,stucki2023,lux2025}. Este trabajo
estudia esas herramientas y las pone a prueba en un problema real.

El problema elegido es la segmentación de filamentos solares en
imágenes H-alfa, propuesto en la competición \emph{Solar Filament
Segmentation Challenge 2026}~\cite{kaggle2026}. Los filamentos son
estructuras alargadas, con ramificaciones finas llamadas \emph{barbas},
y la competición los evalúa con una métrica de instancias, la calidad
panóptica~\cite{kirillov2019}, que castiga con fuerza los filamentos
partidos o fusionados. Es, por tanto, un caso en el que los errores
topológicos tienen un coste directo y medible.

% ------------------------------------------------------------
\section{Objetivos}
% ------------------------------------------------------------

El objetivo general es estudiar qué aporta la topología computacional a
la segmentación de estructuras finas, tanto en la teoría como en la
práctica. Se concreta en cuatro objetivos:
\begin{enumerate}
  \item[\textbf{O1.}] Estudiar los fundamentos de la homología y la
  homología persistente, y los métodos de segmentación que se basan en
  ellas, en particular Topograph~\cite{lux2025}.
  \item[\textbf{O2.}] Construir un protocolo de experimentación
  reproducible para la segmentación de filamentos, en el que las
  comparaciones entre métodos no dependan de decisiones tomadas sobre
  los datos de evaluación.
  \item[\textbf{O3.}] Diagnosticar los errores de un modelo de
  referencia e interpretarlos en términos topológicos.
  \item[\textbf{O4.}] Evaluar si las pérdidas con sesgo topológico
  (clDice y Topograph) mejoran a las pérdidas de solapamiento, en la
  métrica de la competición y en la detección de las barbas, y cómo
  interactúan con el post-proceso.
\end{enumerate}

% ------------------------------------------------------------
\section{Contribuciones}
% ------------------------------------------------------------

\begin{itemize}
  \item Una exposición de la homología persistente y de Topograph sin
  demostraciones, centrada en la intuición y con ejemplos, pensada para
  un lector con formación matemática pero sin experiencia previa en
  topología computacional.
  \item Una lectura topológica del post-proceso habitual (umbral y
  componentes conexas): el número de instancias predichas es el número
  de puntos de un cuadrante del diagrama de persistencia $H_0$, y cada
  tipo de error deja en él una firma reconocible.
  \item Una comparación controlada de diez pérdidas, combinando BCE con
  Dice, IoU y clDice, con y sin Topograph, bajo un mismo protocolo, con
  separación entre ajuste y reporte e intervalos de confianza.
  \item El resultado principal: Topograph mejora a la baseline sin
  post-proceso, sobre todo porque elimina componentes espurias. Pero el
  filtro de área corrige ese mismo tipo de error, y con él la ventaja
  desaparece.
  \item Una medida de la detección de las barbas basada en la espina
  anotada en MAGFiLO.
  \item El código, público, con el que se puede reproducir cada
  resultado.
\end{itemize}

% ------------------------------------------------------------
\section{Estructura de la memoria}
% ------------------------------------------------------------

El Capítulo~\ref{cap:fundamentos} presenta los fundamentos de
topología computacional y los métodos topológicos de segmentación. El
Capítulo~\ref{cap:contexto} describe el problema: los filamentos
solares, el conjunto de datos MAGFiLO, la competición y los trabajos
previos. El Capítulo~\ref{cap:planteamiento} diagnostica los errores
de un modelo de referencia, los interpreta topológicamente y plantea la
estrategia de mejora. El Capítulo~\ref{cap:metodologia} detalla el
protocolo experimental, y el Capítulo~\ref{cap:resultados}, los
resultados. El Capítulo~\ref{cap:conclusiones} recoge las conclusiones,
las limitaciones y el trabajo futuro.
